% Source: https://artofproblemsolving.com/community/c3h2902301_whats_wrong_with_this_solution
% Stable identity: aops-topic-2902301-q01
Substituting $x^2=y+a$: $y^2+y+a-a^2=0$, i.e. $(y+a)(y-a+1)=0$. The root $y=-a$ gives $x=0$, a single point. The root $y=a-1$ gives $x^2=2a-1$, which yields two more points exactly when $2a-1>0$. For example, $a=1$:
\begin{asy}
usepackage("amsmath");
size(3cm);
defaultpen(fontsize(10pt));
import graph;draw(circle((0,0),1));draw(graph(new real(real x){return x^2-1;},-1.6,1.6));dot((0,-1));dot((1,0));dot((-1,0));draw((-1.8,0)--(1.8,0),gray,Arrows(4));draw((0,-1.4)--(0,1.8),gray,Arrows(4));
\end{asy}
 So there are exactly three points iff $\boxed{a>\dfrac12}$.

The flaw: requiring distinct $y$-roots is not enough, since each $y$-root yields $0$, $1$ or $2$ points depending on the sign of $x^2=y+a$. One must count points, not roots in $y$.
